\documentclass[11pt]{article}

\usepackage[margin=25mm]{geometry}

\usepackage{amsmath,amssymb,fontspec,unicode-math}

\IfFontExistsTF{DejaVu Serif}{\setmainfont{DejaVu Serif}}{\setmainfont{Latin Modern Roman}}

\IfFontExistsTF{Latin Modern Math}{\setmathfont{Latin Modern Math}}{\setmathfont{DejaVu Math TeX Gyre}}

\usepackage{xurl,hyperref}

\setlength{\parindent}{0pt}\setlength{\parskip}{6pt}

\begin{document}

\title{BeyondMix: Leveraging Structural Priors and Long-Range Dependencies for Domain-Invariant LiDAR Segmentation}

\author{ELVIA preliminary AI analysis}\date{21 October 2026}\maketitle

Preliminary assessment: Minor.\par

Independent scientific review remains pending. Author adoption has not been requested.

Cross-domain receptive fields are neither universally necessary nor, by the stated proof alone, sufficient for invariant useful features. These limitations concern the supporting distributional argument rather than the empirical segmentation method.

\section*{Review scope}

Latest-version, claim-focused review of the cited author-paper passages and applicable current author-release paths. The stated counterexamples and scope controls were checked analytically; no author program, old verifier, training run or benchmark was executed.

\section*{Problem 1: A local feature can remain invariant and useful under both shifts}

Source: Appendix A, PDF p.21, reverse proposition, Eqs.(18)–(20); Section 1, PDF p.3

Assessment: supported; impact: minor

\subsection*{Source claim (paraphrase)}

Appendix A’s reverse implication rules out effective invariant features without a receptive field spanning both domains.

\subsection*{Analysis}

Let \(X=(U,A)\), with balanced \(U\), nuisance probabilities \(P_s(A=1)=1/4\), \(P_t(A=1)=3/4\), and \(P_s(Y=U\mid A)=0.95-0.20A\), \(P_t(Y=U\mid A)=0.75+0.20A\). Both raw and conditional laws shift. Yet the local feature \(f(X)=U\) has identical balanced domain laws, zero domain information and zero population MMD; its Bayes classification accuracy is \(0.90\) in each domain.

\subsection*{Finding and scope}

The universal necessity implication fails. The usefulness of larger receptive fields on the evaluated data remains possible.

\subsection*{Proposed correction}

State a distribution-, hypothesis-class- and optimization-dependent claim rather than universal necessity. No change to the empirical method follows from this scope correction.

An explicit local correction is described; author adoption and a corresponding published erratum are not established.

\section*{Problem 2: Mixed-input gradients do not follow from the displayed weighted objective}

Source: Appendix A, PDF p.21, forward argument, Eqs.(15)–(17)

Assessment: supported; impact: minor

\subsection*{Source claim (paraphrase)}

Appendix A’s forward argument identifies a mixed-input gradient with a weighted source/alignment gradient and then asserts near-zero domain information.

\subsection*{Analysis}

For inputs \(x_s=1\), \(x_t=3\), mixture weight \(1/2\) and a logistic linear feature at parameter zero, the mixed supervised gradient is \(-1\), while the displayed weighted gradient is \(-1/4\). If Eq.(16) defines a weighted objective instead, the algebra is valid but concerns another object. Equation (17) additionally assumes the near-equality \(H(D\mid Z)\approx H(D)\) needed for invariance.

\subsection*{Finding and scope}

The written argument does not establish sufficiency for the actual learning process. It does not prove that a specific optimizer fails.

\subsection*{Proposed correction}

Specify the mixing, labels, feature class, loss and optimization conclusion. Under a characteristic kernel, zero population MMD implies equal feature laws, but receptive-field coverage alone does not establish that premise.

An explicit local correction is described; author adoption and a corresponding published erratum are not established.

\section*{Limits}

Preliminary AI assessment; independent human scientific adjudication remains pending.\par

An analytic counterexample or a release discrepancy does not by itself remeasure the paper’s reported experiments.\par

\section*{Source identity}

NeurIPS 2025 official proceedings PDF (current publisher bytes, checked 2026-10-21)

\url{https://proceedings.neurips.cc/paper_files/paper/2025/file/812d5b873efe714cfc26342436ea0b0d-Paper-Conference.pdf}

Reviewed PDF SHA-256: \nolinkurl{92b101b1cff89497e2ca9ba65b94ca96b2d9a741850df830f1534ddb868ecad1}

The current official publisher PDF is the same work as the original case: title, author identity and publisher hash URL agree, and the freshly downloaded scientific PDF is byte-identical to the prior source. The actual download timestamp and SHA-256 are recorded for Issue \#27. Exact-title arXiv search did not return a matching later version; background references were excluded.

\end{document}
